\sscp{Preposed possessors and alienable-inalienable distinctions}{02-03-01-01}
Perhaps the most well known feature of the Austronesian
languages of Wallacea is the \it{preposed possessor} (also known as the ``preposed genitive''
or the ``reverse genitive'')
This construction was introduced briefly in \srf{01-05}.
It is not found in the western
and southern parts of the \tsc{\ili{Bima}-Lembata} subgroup,
although it is widespread in other subgroups in our region.
The preposed possessor is also found west of the Blust line in
\ili{Banggai}, setting that language apart typologically from other \tsc{\ili{Saluan}-Banggai}
languages that share a common historical phonology \citep[68]{Mead2003b}.
\ili{As} mentioned previously, prepositional SVO languages with preposed
possessors are typologically unusual in the world’s languages
(see \crf{ch:IntCom}, footnote \ref{fn:01-03}), yet that pattern represents
the majority of Austronesian
languages in the target region.
This typological aberration is not readily explained
from within the languages themselves.
Structural change due to intense contact over a long period of
time with other languages that already have the preposed possessor
pattern common to their own SOV typology provides a far more
credible explanation.\footnote{
		In a worldwide sample,
		preposed possessors are found in only 6--9{\%} of SVO languages
		that have prepositions \citep{Dryer2013a,Dryer2013b,Dryer2013d}.
		By contrast, \citet[118]{Schapper2015} notes, “almost all Papuan languages
		have the noun following the genitive (95{\%} in the sample of
		\citealp{Dryer2013b} [\ldots]”,
		referring to the pattern of preposed possessors.
		This fact is borne out by the even more detailed survey
		of \citet[276--285]{SchapperGasser2023}.}

Also in the realm of possession,
while \citet[259]{Blust1993} wanted to ascribe the development of a
distinction between \it{alienable} and \it{inalienable possession} 
in Austronesian languages \ili{east} of
the Blust line to a morphosyntactic innovation in support of
PCEMP, other scholars
\citep{KlamerReesinkVanStaden2008,Schapper2015,ReesinkDunn2018,
Klamer2019,SchapperGasser2023,Grimes2024}
show evidence to attribute the alienable-inalienable
distinction to contact with Papuan languages.
This considerably weakens its value as a subgrouping argument.

In reviewing the literature on possession in the region,
it is clear there is some confusion of both terminology and categories.
Different investigators often use the same term to refer to different
issues or structures associated with possession.
Here we clarify some of the issues that we have found helpful
to understand various discussions:

\begin{itemize}
	\item Broadly speaking,
				two different kinds of possessive constructions need to be distinguished in grammatical descriptions.
				Predicative possession is clausal
				and tends to assert ownership,
				but can also assert a close association (\it{I have/own/possess a house/father/arm.
				/ The machete is mine.
				/ It is mine.}).
				Many language sketches do not discuss or illustrate predicative possession.
				Phrasal possession (adnominal possession) is found in the grammars
				as a type of noun phrase.
				In many languages, only nouns or derived (abstract) nouns can
				be possessed in this way (\it{pig’s head,
				bird’s wing, bird’s nest,
				house roof, house post,
				tree root, my house/father/arm; my journey,
				my return; my faith; my love,
				my hope, my disappointment, my return fare}).
				Some investigators note that the language they are describing
				has both pre-posed and post-posed possessors, whereas on inspection of the actual
				data, the difference is between phrasal possession and predicative possession.
	\item Many word lists tend to illustrate only parts of some
				phrasal possession constructions without enough information to
				see the order of \tsc{possessor}
				and \tsc{possessed} with explicit possessors.
				The presence of possessive clitics in word lists does not actually
				tell us the order of possessor-possessed when the possessor is explicit.
				So abstractions of words of common Wallacean forms that we might
				find in word lists like \it{limana} /lima=na/ ‘hand-\tsc{3sg.gen}’,
				\it{uluŋu} /ulu=ŋu/ ‘head-\tsc{1sg.gen}’, \it{walin} /wali=n/
				‘younger sibling (same sex)-\tsc{3sg.gen}’ do not tell us the
				position of an explicit possessor in those languages.
				They are not post-posed per se.
				We need to find fuller phrasal examples with a pronoun or possessor
				noun to know where the possessor occurs,
				such as \it{\tbf{nia} lima=n} ‘\tsc{3sg} hand-\tsc{3sg.gen}’,
				\it{\tbf{aku} ulu=ŋu} ‘\tsc{1sg} head-\tsc{1sg.gen}’,
				\it{\tbf{niŋ} wali=n} ‘\tsc{3sg.poss} younger sibling (same sex)-\tsc{3sg.gen}’,
				\it{\tbf{oo} ni luma} ‘\tsc{2sg poss} house’.
				In these cases we can see that the possessor,
				when explicit, is pre-posed in all cases.
				And this often matches the order of NOUN + NOUN possessive constructions
				such as \it{\tbf{hahi} ulu=n} ‘pig head=\tsc{3sg.gen}’,
				also with a pre-posed possessor.
				In the \tsc{\ili{Ambon}-Seram} languages that use enclitics to mark
				inalienable possession, and proclitics to mark alienable possession,
				nevertheless when there is an explicit possessor it is pre-posed
				in both constructions (see \trf{02-01}).
				However, with many Papuan languages of TAP
				and \tsc{North Halmahera}, the preposed possessor proclitics/prefixes
				are often the only possessor to be found.
				We also note that in natural texts in most languages,
				an explicit possessor is often omitted from possessive constructions
				because it is known from the discourse context.
				Sometimes explicit possessors are most commonly used in natural
				texts to disambiguate \tsc{3sg} from \tsc{3pl},
				or \tsc{1pl.excl} from \tsc{2pl},
				depending on where ambiguity is found in a particular language.
	\item The term “possessive construction” is slightly misleading
				in that only occasionally do “possessive constructions” illustrate
				what might be “owned”,
				“possessed” or “controlled” in a legal sort of understanding.
				More often various ``possessive constructions'' in the region are
				used to illustrate physical \it{part-whole} relationships (body
				parts, house parts, tree parts,
				boat parts), conceptual part-whole relationships or close associations
				such as many (close) kin relationships (\it{my father,
				my opposite sex sibling}),
				positional relations (\it{top of house,
				rear of boat, edge of garden}),
				or abstract nouns (\it{her story,
				his liking, her hope, his understanding}).
	\item There can be ambiguity as to the particular nature of
				the ``ownership'' or association.
				For example, \ili{Buru} \it{nak endohi-n} ‘his/her story’,
				or \ili{Kotos} \ili{Amarasi} \it{in reta-n} ‘his/her story’ are both ambiguous
				between ‘the story s/he told’ versus ‘the story told about him/her’.
				The \ili{Buru} term can also have the sense of ‘the kin group origin
				story that he has the rights to tell
				and others don’t, but he will only tell it under the right circumstances
				for fear of punishment from the ancestors’.
				It is a story he does not tell,
				but does have the right to tell.
	\item The term “genitive construction” is unhelpful for
				some, and useful to others.
				Those who find it unhelpful see the grammar of the so-called
				“genitive” constructions in the region as very different
				from the genitive case described for Indo-European and classical languages.
				In Indo-European languages, the genitive case is marked on the
				\tsc{possessor} (\it{pig’\tbf{s} head}), whereas in Austronesian languages in the region,
				the “genitive” marker is indexed on the \tsc{possessed} (\it{fafu olo-\tbf{n}} ‘pig head-\tsc{3sg.gen}’).
				While the marking is on different constituents,
				nevertheless the semantic relationship is exactly the same as
				in ‘\it{pig’s head}’ or ‘\it{head of pig}’.
				Furthermore, Indo-European languages tend to be predominantly
				dependent-marking, whereas the Austronesian languages in our target region tend
				to be predominantly head-marking,
				and consequently the genitive marking indexed on the \tsc{possessor}
				(dependent) in Indo-European languages,
				but on the \tsc{possessed} (head) in Austronesian languages in
				the region is entirely consistent with that distinction (see
				\citealp{Nichols1986,NicholsBickel2013}).
				Some investigators find the terms “genitive”
				and “possessive” useful for distinguishing two different grammatical
				constructions of phrasal possession (see below).
	\item Within the category of phrasal (adnominal) possession,
				many investigators make a further distinction between \it{alienable}
				(general, indirect) possession, and \it{inalienable} (obligatory, direct) possession.
				However, \citet{GrimesCIP} points out the discussion of \it{alienable}
				and \it{inalienable} possession in many language descriptions
				in the region confuses the distinction between semantic issues and morphosyntactic issues.
				He suggests the terms \it{alienable} (indirect)
				and \it{inalienable} (direct) be reserved for discussion of the
				semantics and noun class systems.
				And that different possessive constructions in the grammar should
				use different terms, such as a general \it{possessive} construction,
				and a \it{genitive} construction (also an \it{edible} construction
				for \ili{Selaru}, and others as appropriate).
				The reason for making this distinction is not only to clarify
				the confusion in the discussion (where “inalienable” sometimes
				refers to a noun class
				and other times to the grammatical construction often associated
				with it), but also because it appears that while all languages in the region
				seem to have at l\ili{east} two separate constructions in the grammar
				dealing with phrasal possession,
				not all languages have a semantic or noun class distinction of alienable versus inalienable.
				In some languages, \it{all} nouns,
				whether they are considered alienable or inalienable in other
				languages, can be used with the general (indirect) possessive construction.
				And terms that are elsewhere considered inalienable,
				can be used with or without the genitive enclitic as a bare root,
				so it cannot be said to be a bound root,
				and cannot be said to be “inalienable”.
\end{itemize}

We turn first to the evidence for contact-induced change.
\ili{Malay}, west of the Blust line
and outside the target region,
has the typical western pattern with post-posed possessors.
Standard \ili{Malay} (as represented by Indonesian) also does not distinguish
alienable and inalienable noun classes,
nor does it have two phrasal possessive constructions,
as shown in example \qf{02-01}.

{\wg\tblr{>{\hspace{-\tabcolsep}}l<{\hspace{-\tabcolsep}}>{\hspace{-\tabcolsep}\enspace}lll}
\lsptoprule
\tbx{02-01}	&		&	Indonesian	&	gloss\\ \midrule
	&	a.	&	rumah saya	&	‘my house’	\\
	&		&	rumah-ku	&	‘my house’	\\
	&		&	rumah kita	&	‘our (incl.) house’	\\
	&		&	rumah-ɲa	&	‘her/his house’	\\
	&	b.	&	kaki saya	&	‘my foot’	\\
	&		&	kaki-ku	&	‘my foot’	\\
	&		&	kaki kita	&	‘our (incl.) foot’	\\
	&		&	kaki-ɲa	&	‘her/his foot’	\\
	&	c.	&	ayah saya	&	‘my father’	\\
	&		&	ayah-ku	&	‘my father’	\\
	&		&	ayah kita	&	‘our (incl) father’	\\
	&		&	ayah-ɲa	&	‘her/his father’	\\
\lspbottomrule
\end{tabular}\mbox{}\\}

Regional Malays (not formal Indonesian) have another possessive
construction with a possessive word \it{puɲa,
puŋ, pu}, that can be used both predicatively
and phrasally, as seen in example \qf{02-02} using \ili{Kupang} \ili{Malay}.
Notice the possessor in these eastern Indonesian \ili{Malay} constructions
are pre-posed.

{\wg\tblr{>{\hspace{-\tabcolsep}}l<{\hspace{-\tabcolsep}}>{\hspace{-\tabcolsep}\enspace}ll}
\lsptoprule
\tbx{02-02}	&	\ili{Kupang} \ili{Malay}	&	gloss	\\	\midrule
	&	beta puŋ ruma	&	‘my house’	\\
	&	beta ada puŋ ruma.	&	‘I have/own a house.’ (\it{ada} ‘exist, have attribute’)	\\
	&	ada beta puŋ ruma.	&	‘I have/own a house.’ 	\\ \hhline{~}
	& & (pragmatically emphatic, insistent)\\
	&	beta puŋ ruma, ada.	&	‘I have/own a house.’ (lit: my house, exists)	\\
	&	beta puŋ kaki	&	‘my feet’	\\
	&	beta ada puŋ kaki.	&	‘I have feet.’	\\
	&	beta puŋ papa	&	‘my father’	\\
	&	beta ada puŋ papa.	&	‘I have a father.’ (= my father is still living)	\\
\lspbottomrule
\end{tabular}\mbox{}\\}

\largerpage
In contrast to the western pattern of post-posed possessors,
SOV Papuan languages in the target region tend to have pre-posed
possessors and a distinction between alienable
and inalienable noun classes,
each type using a different possessive construction \citep{Klamer2017}.
This is illustrated in examples \qf{02-03}--\qf{02-05} from \ili{Wersing},
a TAP language spoken in eastern \ili{Alor} (data from \citealp{Jones2023}).
The pronominal possessor is prefixed to inalienable nouns in example \qf{02-03}.
We have found no examples in \ili{Wersing} texts where the possessor
of an inalienable noun is further distinguished by a free pronoun.

{\wg\tblr{>{\hspace{-\tabcolsep}}l<{\hspace{-\tabcolsep}}>{\hspace{-\tabcolsep}\enspace}ll}
\lsptoprule
\tbx{02-03}	&	\ili{Wersing}	&	gloss	\\ \midrule
	&	n-teŋ	&	‘my hand’		\\
	&	a-teŋ	&	‘your(sg.) hand’		\\
	&	g-teŋ	&	‘his/her hand, their hands’		\\
	&	t-teŋ	&	‘our (incl.) hands’		\\
	&	ɲ-teŋ	&	‘our (excl.) hands’		\\
	&	y-teŋ	&	‘your (pl.) hands’		\\
\lspbottomrule
\end{tabular}\mbox{}\\}

A different CV-shaped pronominal possessor is preposed to indicate
general (indirect, alienable) possession, as in \qf{02-04}.
Note that the consonants distinguishing person
and number remain the same for all affixes
and pronouns, regardless of whether they are modifying nouns
or verbs in \ili{Wersing}.

{\wg\tblr{>{\hspace{-\tabcolsep}}l<{\hspace{-\tabcolsep}}>{\hspace{-\tabcolsep}\enspace}ll}
\lsptoprule
\tbx{02-04}	&	\ili{Wersing}	&	gloss	\\ \midrule
	&	ne spir	&	‘my boat’		\\
	&	ge spir	&	‘his boat, their boat’		\\
	&	ne sekok	&	‘my chicken’		\\
	&	ge sekok	&	‘his chicken, their chicken’		\\
\lspbottomrule
\end{tabular}\mbox{}\\}

Many body parts in \ili{Wersing} use the general possession construction,
as shown in \qf{02-05}.
So the distinction is not one of part-whole,
or direct versus indirect possession,
but of different body part terms being in different noun classes.

{\wg\tblr{>{\hspace{-\tabcolsep}}l<{\hspace{-\tabcolsep}}>{\hspace{-\tabcolsep}\enspace}lll}
\lsptoprule
\tbx{02-05}	&	\ili{Wersing}	&	gloss	\\ \midrule
	&	ne lkomaŋ	&	‘my head’		\\
	&	e lkomaŋ	&	‘your (sg.) head’		\\
	&	ge lkomaŋ	&	‘his/her head, their head’		\\
	&	te lkomaŋ	&	‘our (incl.) head’		\\
	&	ge lkomaŋ du	&	‘their heads’		\\
	&	tenuku te lkomaŋ	&	‘heads of the two of us (incl.)’ 		\\
	&	ninuku ni lkomaŋ	&	‘heads of the two of us (excl.)’		\\
	&	yenuku ye lkomaŋ	&	‘heads of the two of you (pl.)’		\\
\lspbottomrule
\end{tabular}\mbox{}\\}

With NOUN + NOUN part-whole relationships in \ili{Wersing},
the \tsc{possessor} also precedes the \tsc{possessed},
as in \qf{02-06} and \qf{02-07}.
The possessive free pronoun with an alienable noun in \qf{02-06} is
optional, whereas in \qf{02-07} the possessive prefix marked on the phrasal head
with an inalienable noun is obligatory.

{\wg\tblr{>{\hspace{-\tabcolsep}}l<{\hspace{-\tabcolsep}}>{\hspace{-\tabcolsep}\enspace}lll}
\lsptoprule
						&	\ili{Wersing}	&	gloss	\\ \midrule
\tbx{02-06}	&	\tbf{pei} lkomaŋ	&	‘pig’s head’	\\
						&	\tbf{pei} ge lkomaŋ	&	‘pig’s head’	\\ \midrule
\tbx{02-07}	&	\tbf{kutual} g-teŋ	&	‘dog’s (front) paw’	\\
\lspbottomrule
\end{tabular}\mbox{}\\}

With that background of the western pattern having post-posed
possessors and the Papuan pattern of pre-posed possessors with an additional
alienable-inalienable distinction in noun classes
and grammatical constructions, we turn now to explore some of
the variety and complexity of possessive constructions in the Austronesian
languages in the target region.

In some Austronesian languages there is a strong correlation
between \it{inalienable} nouns with the genitive construction
(marked with genitive enclitics on the phrasal head,
the \tsc{possessed}), and \it{alienable} nouns with the general
possessive construction (often marked with possessive words or possessive proclitics).
But most languages allow some crossover between the two,
and some allow a lot of crossover.

While \citet{LaidigLaidig1990} find an alienable-inalienable
distinction in \ili{Larike} nouns (an \tsc{\ili{Ambon}-Seram} language),
they also find a more complex skewing between form
and function that requires greater nuancing,
“Preliminary data from \ili{Larike} suggest that the notion of bound
nouns (as opposed to inalienably possessed nouns) may provide
a more accurate description” (1990:107).

\citet[xxviii--xxix]{Collins2003b} describes the noun class distinction
for \ili{Asilulu} (an \tsc{\ili{Ambon}-Seram} language)
and illustrates the two different grammatical constructions associated
with them, both with preposed possessors shown in \trf{02-01}
He also notes, “In general,
there is no crossover from one category to another.
A few exceptions to this generalization include loanword kinship
terms, usually classed in the alienable category,
but occasionally appearing with inalienable marking” (p.xxvi).
The occurrence of the “full emphatic pronouns” is optional in
what we call the genitive construction associated with inalienable nouns.
Note that where they occur,
all overt possessors are preposed in \ili{Asilulu}.
Note also that the distinction in two different grammatical constructions
is more clearly seen with singular pronouns,
and less evident with plural pronouns.

\begin{table}[h]
	\centering\tcp{Alienable-inalienable distinction in Asilulu}{02-01}
	\st{5.5pt}\begin{tabular}{lll|lll}\lsptoprule
\mc{2}{l}{Alienable N} & Gloss & \mc{2}{l}{Inalienable nouns} & Gloss\\
\mc{3}{l|}{possessive, indirect}	& \mc{2}{l}{genitive, direct} & \\ \midrule 
aʔu	&	\tbf{ku}-haka	&	‘my boat’	&	(aʔu)	&	ulu-\tbf{ku}	&	‘my head’		\\
ale	&	\tbf{mu}-buku	&	‘your (sg.) book’	&	(ale)	&	wali-\tbf{mu}	&	‘your (sg.) ySib’		\\
ali	&	\tbf{na}-luma	&	‘his house’	&	(ali)	&	nala-\tbf{ni}	&	‘his name’		\\
ami	&	\tbf{ma}-tipil	&	‘our (excl.) basket’	&	(ami)	&	\tbf{ma}-tuku	&	‘our (excl.) knees’		\\
ite	&	\tbf{ra}-pikal	&	‘our (incl.) plate’	&	(ite)	&	mata	&	‘our(incl.) eyes’		\\
imi	&	\tbf{mi}-lapun	&	‘your (pl.) shirt’	&	(imi)	&	\tbf{mi}-meme	&	‘your (pl.) MB’		\\
sini	&	\tbf{ri}-kata	&	‘their pants’	&	(sini)	&	\tbf{ri}-lima	&	‘their arms’		\\
		\lspbottomrule
	\end{tabular}
\end{table}

The nearby \tsc{Sula-Buru} language of \ili{Buru} works very differently from \ili{Asilulu}.
\citet[85--86, 190--192, 277--292, 330--332]{Grimes1991c} describes
a number of structural,
semantic, and typological issues associated with different types
of “possession” in \ili{Buru},
all of which have preposed possessors.
He notes that \it{all} nouns,
regardless of whether they are bound or free,
can be used with the general \it{possessive} construction.
And \it{any} noun that can be perceived in a physical or conceptual
part-whole relationship can be used in the \it{genitive} construction.
Many nouns can use both constructions at the same time.
Examples \qf{02-08}--\qf{02-09} show the \ili{Buru} general possessive construction
with both a so-called “alienable” noun (\it{huma} ‘house’),
and an “inalienable” body part (\it{olo} ‘head’).
Examples \qf{02-10}--\qf{02-11} show a body part using the \ili{Buru} genitive construction.
All possessors are preposed.

\begin{exe}
\setcounter{xnumi}{7}
	\ex{\gll	naŋ huma\\
						\tsc{1sg.poss} house\\
			\glt	`my house' (\it{huma} < *Rumaq)
						}\label{ex:02-08}
	\ex{\gll	naŋ olo\\
						\tsc{1sg.poss} head\\
			\glt	`my head' (\it{olo} < *qulu)
						}\label{ex:02-09}
	\ex{\gll	(ya) olo=ŋ\\
						\tsc{1sg} head=\tsc{1sg.gen}\\
			\glt	`my head'
						}\label{ex:02-10}
	\ex{\gll	(fafu) olo=n\\
						pig head=\tsc{3sg.gen}\\
			\glt	`pig's head'
						}\label{ex:02-11}
\end{exe}

In \ili{Buru}, body parts tend to be associated with a whole,
and therefore are commonly found with the genitive enclitics,
as in example \qf{02-12}.

\begin{exe}
	\ex{\gll	da iko tu \tbf{olo=n}\\
						\tsc{3sg} go with head=\tsc{3sg.gen}\\}\jambox*{(\tsc{genitive} -- \tsc{no possessive})}
			\glt	‘He went with its (pig’s) head.’
						(accompanied possession that may belong to someone else)
						\label{ex:02-12}
\end{exe}

When both the possessive construction
and the genitive construction are used together in \ili{Buru},
each is mapping a different semantic relationship,
as in example \qf{02-13}.
The genitive construction indicates the head is part of a whole
(pig); the possessive construction indicates it is owned or in
the possession of the subject of the clause.

\begin{exe}
	\ex{\gll	da iko tu \tbf{nak} \tbf{olo=n}\\
						\tsc{3sg} go with \tsc{3sg.poss} head=\tsc{3sg.gen}\\}\jambox*{(both \tsc{possessive} and \tsc{genitive})}
			\glt	‘He went with his (pig’s) head.’
						(accompanied possession that he owns)
						\label{ex:02-13}
\end{exe}

Other recorded text examples also show that physical body parts
in \ili{Buru} do not always require genitive enclitics,
particularly when followed by attributive modifiers, as in \qf{02-14}.
Therefore they are not bound nouns.

\begin{exe}
	\ex{\gll	Susana \tbf{nak} \tbf{olo} em-hapu\\
						Susana \tsc{3sg} head \tsc{stat}-tie\\}\jambox*{(\tsc{possessive} -- no \tsc{genitive})}
			\glt	‘Susana had her (own) head tied up (in a cloth).’
						%(Showing her working hard.)
						\label{ex:02-14}
\end{exe}

In \tsc{Sula} (also \tsc{Sula-Buru}),
all nouns can be possessed by juxtaposing the possessor before
the possessed \citep[280--293]{Bloyd2020},
and there is also an uninflected possessive word that is used
in some third person combinations.

The \tsc{\ili{Bima}-Lembata} language \ili{Dhao} \citep{Grimes2010a} works
quite differently from both \ili{Asilulu}
and \ili{Buru} regarding possession.
There is no structural distinction between what in other languages
might be considered alienable or inalienable nouns.
Both types have post-posed possessors (free pronouns or enclitics)
as shown in \trf{02-02}.
Note that in \ili{Dhao} penultimate schwa is syllabic
and stress bearing, and triggers phonetic lengthening of the following adjacent consonant.

\begin{table}[h]
	\centering\tcp{Postposed possessors in Dhao}{02-02}\begin{threeparttable}
	\st{4.5pt}\begin{tabular}{lll|lll}\lsptoprule
\mc{2}{l}{`house'} & Gloss & \mc{2}{l}{`head'\tnote{†}} & Gloss \\ \midrule
əmu	&	\tbf{ʤaʔa}	&	house \tsc{1sg} 	&	kətu	&	\tbf{ʤaʔa}	&	head \tsc{1sg} 		\\
	&		&	‘my house’	&		&		&	 ‘my head’		\\
əmu	&	\tbf{əu} {\tl} \tbf{mu}	&	house \tsc{2sg}	&	kətu	&	\tbf{əu} {\tl} \tbf{mu}	&	head \tsc{2sg}		\\
əmu	&	\tbf{nəŋu} {\tl} \tbf{na}	&	house \tsc{3sg}	&	kətu	&	\tbf{nəŋu} {\tl} \tbf{na}	&	head \tsc{3sg}		\\
əmu	&	\tbf{ʤiʔi}	&	house \tsc{1pl.excl}	&	kətu	&	\tbf{ʤiʔi se}	&	head \tsc{1pl.excl}		\\
əmu	&	\tbf{ə}{\tbfDZ}\tbf{i}	&	house \tsc{1pl.incl}	&	kətu	&	\tbf{ə}{\tbfDZ}\tbf{i se}	&	head \tsc{1pl.incl}		\\
əmu	&	\tbf{miu} {\tl} \tbf{mi}	&	house \tsc{2pl}	&	kətu	&	\tbf{miu se}	&	head \tsc{2pl}		\\
əmu	&	\tbf{rəŋu} {\tl} \tbf{ra}	&	house \tsc{3pl} 	&	kətu	&	\tbf{rəŋu} {\tl} \tbf{ra səra}	&	head \tsc{3pl} \\
		&	&	‘their house (sg.)’	&	&	&	‘their heads (pl.)’		\\
		\lspbottomrule
	\end{tabular}
			\begin{tablenotes}\tnsize
				\item[†]	{
									Without the plural deictics (\it{se, səra})
									with the plural pronouns in \ili{Dhao},
									these would be interpreted as a single social head ‘their boss’,
									rather than plural physical heads of multiple bodies ‘their heads’.}
			\end{tablenotes}
	\end{threeparttable}
\end{table}

But like other languages in the region,
\ili{Dhao} has at l\ili{east} two ways of expressing possession,
including with a possessive word \it{unu} ‘own, possess, have’.
The possessor is still post-posed,
as illustrated in examples \qf{02-15}--\qf{02-17}.
The \it{unu} word can be used both predicatively and phrasally.
Used phrasally (adnominal possession),
it is often optional,
often being used to break up a heavy information load without
much difference in pragmatic force.

\begin{exe}
	\ex{\gll	nəŋu pa-taleo rai na / nəŋu\\
						\tsc{3sg} \tsc{caus}-wide land \tsc{3sg.gen} / \tsc{3sg}\\
			\glt	`He enlarged his land.'
						}\label{ex:02-15}
	\ex{\gll	nəŋu pa-taleo rai \tbf{unu} na / nəŋu\\
						\tsc{3sg} \tsc{caus}-wide land own \tsc{3sg.gen} / \tsc{3sg}\\
			\glt	`He enlarged his land.' / ‘He enlarged the land owned by him.’
						}\label{ex:02-16}
	\ex{\gll	nəŋu pa-taleo rai \tbf{unu} Sarus\\
						\tsc{3sg} \tsc{caus}-wide land own Lazarus\\
			\glt	‘He enlarged the land owned by Lazarus.’
						}\label{ex:02-17}
\end{exe}

The \ili{Dhao} possessive word \it{unu} ‘own,
possess, have’ can also be used to assert possession,
as in example \qf{02-18}.

\begin{exe}
	\ex{\gll	rai əna \tbf{unu} nəŋu\\
						land \tsc{dist} own \tsc{3sg}\\
			\glt	`That land is owned by him.'
						}\label{ex:02-18}
\end{exe}

The \tsc{\ili{Bima}-Lembata} language \ili{Sika} also does not distinguish
between alienable and inalienable nouns \citep[607--608]{LewisGrimes1995}.
With possessive pronouns it further has fluid order
and can mark the same nouns with either a preposed possessor
or a post-posed possessor,
as in the following examples.
The modern genitive suffixes have three different forms
(\it{-n, -ŋ, -t}) and no longer preserve an assumed earlier system of
inflection in the genitive suffixes.
Body parts can be possessed with variable order in the possessive
phrase, as in examples \qf{02-19} and \qf{02-20}.

\begin{exe}
	\ex{\gll	mata \tbf{nimu-n}\\
						eye \tsc{3sg-gen}\\
			\glt	`her/his eye'
						}\label{ex:02-19}
	\ex{\gll	\tbf{nimu-ŋ} mata \\
						\tsc{3sg-gen} eye\\
			\glt	`her/his eye'
						}\label{ex:02-20}
\end{exe}

Proper noun possessors in \ili{Sika} tend to be preposed
and do not directly take a genitive enclitic.
But they also make no distinction between general nouns,
kin terms, and body parts, as in \qf{02-21}.

{
\setcounter{tableex}{20}
\wg\tblr{>{\hspace{-\tabcolsep}}l<{\hspace{-\tabcolsep}}>{\hspace{-\tabcolsep}\enspace}llll}
\lsptoprule
\tbx{02-21}	& Rapa	& muʔu-n	&`Rapa's bananas' &[(name) banana-\tsc{gen}]\\
& Lodan		& nara-t	&`Lodan's brother' &[(name) brother-\tsc{gen}]\\
& Kletus	& baʔi-t	&`Kletus's foot' &[(name) foot-\tsc{gen}]\\
\lspbottomrule
\end{tabular}\mbox{}\\}

With the exception of a few languages like \ili{Dhao},
\ili{Sika} and many other \tsc{\ili{Bima}-Lembata} languages,
typologically both the genitive
and possessive grammatical constructions tend to have preposed
possessors in most languages
and subgroups in the region (a pattern associated with SOV typology),
so the typological abberation with prepositional SVO languages
is still highly significant for those languages,
which along with the presence of an alienable-inalienable distinction
in Papuan languages, supports the explanation of contact-induced change.
But there are also many exceptions
and variables in the region.
So because not all Austronesian languages in several subgroups
in the region have a noun class alienable-inalienable distinction,
and not all languages have preposed possessors,
we need to be cautious about painting the whole region with a
single brush with regard to possession -- both typologically
and genealogically (see \srf{01-05}; \cite[259]{Blust1993}, \cite{Himmelmann2005}).

Looking at the \it{distribution of preposed possessors} in Austronesian
languages in \frf{02-16},
it is clear that they are not found throughout the whole region.
Preposed possessors is the feature first identified nearly 150
years ago that sets Austronesian languages in the Wallacean region
apart from those to the west and north.
The preposed possessor is also one defining feature of \citeauthor{Himmelmann2005}'s (\citeyear{Himmelmann2005})
division of Austronesian languages into two major typological types.

Because of discrepancies of terminology in the descriptions of
languages in the region,
we must clarify what we are representing in \frf{02-16}
and what is not represented.
First of all, we are talking about phrasal (adnominal) possession,
as in \it{my house,
my hand, my father, pig’s head}.
We are not talking about predicative possession,
as in \it{I have a house,
this house is mine}.
The grammar of the two is often quite different.
\frf{02-16} only reflects phrasal possession.

\begin{figure}
	\includegraphics[width=\columnwidth]{possessor.pdf}
	\caption[{Order of the possessor in phrasal possession}]
	{Order of the possessor in phrasal possession\footnotemark}
	\label{fig:02-16}
\end{figure}
\footnotetext{
Data in \frf{02-16} is from \citet{Dryer2013d} supplemented
by our own data sources, as listed in appendix \ref{ch:LanLis}.}

We are mapping phrasal possessive constructions (\it{my father,
pig’s tusk}), and we are excluding vocatives (\it{[my] father!})
indicating direct address, since those often have different grammar,
albeit often having some overlap with possessive systems.
So \ili{Tetun} \it{haʔu-k ama {\tl} haʔu-kan ama} ‘\tsc{1sg-poss} father
= my father’ (possession; preposed),
but \it{ama haʔu-k} ‘father \tsc{1sg-poss} = my father’ (vocative; post-posed).
They are not the same construction.
Or \ili{Buru} \it{(ya) naŋ ama} ‘(\tsc{1sg}) \tsc{1sg.poss} father’,
\it{(ya) naŋ ina} ‘(\tsc{1sg}) \tsc{1sg.poss} mother’ (possessive;
pre-posed), but \it{ama-n} ‘[my/our] father-\tsc{voc}’,
\it{ina-n} ‘[my/our] mother-\tsc{voc}’ (vocative; enclitic).\footnote{
		The \ili{Buru} vocative \it{ama-n} ‘[my/our] father!’ cannot use a free
		pronoun or a preposed possessor/pronoun.
		The possessive construction cannot use an enclitic on \it{ama}:
		\it{naŋ ama} ‘my father’,
		\it{nam ama} ‘your (sg.)/our (incl.) father’,
		\it{nak ama} ‘her father’,
		\it{nan ama} ‘our (incl.) father’,
		\it{nim ama} ‘your (pl.) father’,
		\it{nin ama} ‘their father’ (\it{nun ama} in the \ili{Masarete} dialect).
		But not *\it{naŋ ama=ŋ} ‘my father’,
		nor *\it{ya ama=ŋ} ‘my father’.}
Neighbouring \ili{Lisela} has the
same possessive construction as \ili{Buru},
but a proclitic rather than enclitic in the vocative \it{ŋ-ama}
‘\tsc{1sg.voc}-father’ and \it{ŋ-ina} ‘\tsc{1sg.voc}-mother’.
Vocatives are excluded from \frf{02-16}.

We are mapping the order of the possessor in both NOUN + NOUN
possessive constructions (\it{pig’s head,
roof of house, uncle’s garden,
grandfather’s machete}), and PRONOUN + NOUN possessive constructions
(\it{my house, his hand, their father}).
While most languages have the same order for both (either pre-posed
or post-posed possessors), a few languages have a different order for the possessor in these two constructions.

A few languages have fluid order,
in which the same possessor with the same possessed item can
be found with both orders,
as with \ili{Sika}, as in \it{mata nimu-n {\tl} nimu mata-ŋ} ‘his eye’
\citep[607f]{LewisGrimes1995}.
Nevertheless one order is often dominant,
usually consistent with the pattern in surrounding languages.

Many languages in Wallacea distinguish between alienable and inalienable possession.
This is true for both Papuan
and many Austronesian languages.
But not all of them make this distinction in the morphosyntax.
In some languages the distinction is reflected as a noun class system.
In some it reflects a difference between how bound roots
and free roots are possessed.
In others it is reflected in different grammatical ways of marking possession.
In some there is overlap between all of these parameters.
And in many languages,
different possessive constructions reflect different semantic
relationships, and the different possessive constructions may be used in combination with each other.
To ensure we have checked most common patterns,
we have tried to look minimally at body parts
and other part-whole relationships,
close kin terms, and the possession of general items such as
\it{house, machete, clothes, garden}, etc.

Many languages use genitive/possessive clitics,
regardless of the order of whether the overt possessor is pre-posed or post-posed.
Some languages do not use genitive/possessive clitics,
and all underived nouns are free roots.

We are also excluding topicalised constructions fronted for pragmatic
purposes, or possession in contrastive focus constructions,
since the grammar of those is altered in some languages for pragmatic
focus, and does not always match the normal (unmarked) phrasal possessive construction.
They are marked, and as such they are different.
Categories reflected in \frf{02-16}:

\begin{itemize}
	\item \tbf{Post-posed} possessor (\ili{Malay}: \it{rumah saya} ‘my
				house’, \it{bapak saya} ‘my father’\it{,
				taŋan saya} ‘my hand’,
				\it{kəpala babi} ‘pig’s head’).
	\item \tbf{Pre-posed} possessor (\ili{Buru}: \it{ya naŋ huma} ‘my
				house’, \it{(ya) naŋ ama} ‘my father’,
				\it{yako faha=ŋ {\tl} (ya) naŋ faha=n} ‘my hand/arm’,
				\it{fafu olo=n} ‘pig’s head’).
	\item Both pre-posed
				and post-posed possessors are found in the same language.
				Several different patterns are found:
	\begin{itemize}
		\item	\tbf{Both} orders are found with the same possessor
					and object (\it{my hand {\tl} hand my}).
		\item	The order of NOUN + NOUN possessive constructions is different
					from the order of PRONOUN + NOUN.
		\item	The same pronoun has a different order with different lexical
					items of a similar class such as with different
					body parts or different kin terms.
		\item	There may be a person-based split.
					First person forms have a different order of the possessor from
					second and third person forms.
					Or some other person-based combination.
	\end{itemize}
\end{itemize}

West and north of the Blust line post-posed is dominant,
and this pattern extends \ili{east} of the Blust line through much
of the \tsc{\ili{Bima}-Lembata} subgroup.
\ili{Banggai} is an exception west of the line.
East of \tsc{\ili{Bima}-Lembata} pre-posed is dominant,
with no purely post-posed patterns.
Post-posed possessors are also not found in Papuan languages
of several families in the region for which we have data.
There is a zone of transition from post-posed to pre-posed in
eastern \tsc{\ili{Bima}-Lembata}, specifically in the \tsc{Flores-Lembata} node,
but not in all languages.
And we note that languages that exhibit both orders are limited
to peripheral areas, with some clustered in the \tsc{Flores-Lembata}
node where post-posed
languages meet or transition to pre-posed languages,
and others scattered near the Bird’s Head region of west Papua,
in both Austronesian and Papuan languages.

Pre-posed possessors and an alienable-inalienable distinction
in Austronesian languages in the region are due to intense contact
over a long period of time with Papuan languages that are typologically
SOV and make those distinctions.
\ili{As} \citet{Fricke2019} has pointed out,
the multiple patterns in possessive constructions found within
the \tsc{Flores-Lembata} node of the \tsc{\ili{Bima}-Lembata} subgroup
means that one cannot attribute a pre-posed possessor to the
proto-language of that lower node,
and certainly not to the higher-level subgroup to which it belongs.
The same holds true for the CEMP hypothesis.
The shift from post-posed possessors west of the Blust line to
pre-posed possessors \ili{east} of the Blust line cannot be attributed
to a parent language common to all Wallacean subgroups.